\section{Global scalar-trial accuracy}
\label{sec:global}

The local comparison budget needs an initial error smaller than $x^4$.
We obtain it first for the unmasked field before adding seeds. The pair
inverse upgrades a coarse comparison to this strict saving; only
convergence of the common temperature is needed at the coarse stage.

\begin{proposition}[Accuracy before field seeding]
\label{global:trial-gap}
For the direct paths, use the matrix \(b\) of \eqref{eq:P3} and the
\emph{unmasked, unseeded} field \(h_*=t-bq\) in the untilted system.
Fix \(T>1\) and the exponent choices of Section~\ref{sec:paths}.
For \(m\) in the stated dyad, the following bound is uniform over the
allowed \(\lambda,\ell,\gamma\). There is a saving \(c_G>0\)
independent of the allowed \(p_*\) and of \(s_*\) sufficiently close
to \(1/2\), such that for all sufficiently large \(N\),
\[
 0\le F_*-f_m^0\le x^{4+c_G},\qquad
 F_*= f_{\rm spin}(t)+\tfrac14\int bq^2\,dv .                           \tag{G1}\label{eq:G1}
\]

Here \(f_m^0=f_m(b,h_*)\), and \(f_{\rm spin}(t)\) is the
one-site scalar pressure with the same terminal mass \(\ell\).
The size threshold may depend on the fixed exponent choices.
\end{proposition}

\begin{proof}
Replacing the matrix by the field \(bq\) gives a nonnegative
square-error interpolation, which proves the first inequality.
Throughout the proof normalize \(\ell\) to 1 as in \eqref{eq:A1}:
multiply the pressure by \(\ell\), the clocks by \(\ell^2\), and use
mass \(v/\ell\). The size rescaling is held fixed; below \(\theta\)
denotes a new interpolation parameter. We will remove the matrix in a
constrained interpolation. The next two inputs control its endpoints:
PE bounds the cost of an initial small removal, while scalar coercivity
bounds the endpoint at which the matrix has been removed completely.

\paragraph{The finite-propagation comparison used below}
We use PE, Proposition~6.6, in its ordinary all-mass case
\(I=[0,1]\) (PE, Definition~6.1 and equations~(3)--(6)). Its
parameters are a small fixed power slack \(m_{\rm PE}>0\),
\(\delta_{\rm PE}=m_{\rm PE}/10000\), and scales satisfying
\[
 0<d=O(1),\qquad
 m^{-2/3+m_{\rm PE}}d\lesssim E\le m^{-1/2}d^2,
 \qquad a_{\rm PE}=\sqrt m E/d,\quad v_*=(mE)^{-1}.
\]
Here \(d\) is PE's scale parameter, not our field defect. The bare
clocks \(b_0,h_0\) must be bounded, nonnegative and nondecreasing,
with a fixed positive matrix margin after the comparison; \(d^2\)
has a fixed positive gap below \(T\). At open-interval masses
\(z\gtrsim a_{\rm PE}\), with fixed comparison constants, the
endpoint hypotheses are
\(b_0(z)=T+o(z^2)\) and \(h_0(z)=o(z^3)\).

At strength \(0<\lambda_{\rm PE}\le d^2\), this comparison removes
\(e(r)=\lambda_{\rm PE}f(\lambda_{\rm PE}r/d^3)\) from the matrix,
adds \(re(r)\) to the field, and adds \(\int e(r)r^2/4\) to the
objective. It minimizes over onto quantiles of \([0,1]\) whose
Lipschitz inverses satisfy
\[
\begin{gathered}
 k(r)\le\alpha'(r)\le M_{\rm PE}(\alpha(r)),\qquad
 k(r)=\frac{m^{\delta_{\rm PE}}v_*^2}{(r+m^{2\delta_{\rm PE}}v_*)^2},\\
 M_{\rm PE}(v)=m^{\delta_{\rm PE}}\left(
 \frac{v+v_*}{\min(v_*,a_{\rm PE}^3/d^2)}
 +\frac{d^3}{E}(v+v_*)^2\right).
\end{gathered}
\]
The fixed smooth taper \(f\) equals 1 near zero and is of order
\((1+\log t)/t\) at infinity. Its logarithmic decay slope
\(p(t)=-tf'(t)/f(t)\) increases from 0 towards 1, with
\(1-p\gtrsim(1+\log(1+t))^{-1}\). Provided the changed paths
are admissible, the minimum exceeds the bare pressure by at most
\(m^{20\delta_{\rm PE}}E\), and each minimizing radius is comparable
to mass above \(m^{\delta_{\rm PE}}\sqrt m E/\sqrt{\lambda_{\rm PE}}\).
In our use, the matrix margin and taper monotonicities ensure admissibility.

\paragraph{A common varying base temperature}
The preceding conclusion remains valid when the endpoint hypotheses
use a common \(T_m\to T>1\), with \(T_m\) in a fixed compact subset
of \((1,\infty)\) and held fixed throughout each parent/child tree
at size \(m\). The scale restrictions and propagation slack stay
the same. This sequential extension is proved in
\cite[Appendix: Companion interfaces, paragraph
``A common prelimit temperature'']{PE}.
We recall its argument to specify how it applies here without a rate
of convergence for \(T_m\).

First, the constraints and conditional moment estimates in PE,
Section~3 use bounded clocks and fixed comparability constants;
its endpoint Taylor hypothesis enters only the subsequent profile
argument. At the regular calibration cut in PE, Proposition~4.1,
use \(b_*=T_m+O(W^2)\), \(K_*=T_mS_*+O(W^3)\), and set
\(J_0=T_m\chi^2-1\) in PE~(26)--(27). The proof divides by
\(\chi K_*\asymp W\) and absorbs
\(\eta_n|J_0|x_{\rm test}^2/W^2\). These operations use the actual
\(T_m\), so they introduce no \(T_m-T\) error; the conclusion
PE~(22) likewise uses the actual scalar clock.

In the small-scale proof in PE, Section~5 take time \(K/T_m\), subtract \(K/T_m\) in the scalar compactness and uniform scalar linearization lemmas, and put \(((b_0-T_m)S+h_0)/x_{\rm test}^3\) for the endpoint remainder at test scale \(x_{\rm test}\). All prelimit identities and bounds there are then identical (scalar derivative and strict curvature bounds use only compact temperature and total time bounds); in the rescaled ODE/potential limits \(T_m\) converges to \(T\).

In the macroscopic direct result (PE, Section~5, Proposition~5.9) the actual endpoint convergence and limiting variational problem are at \(T\), so its fixed-temperature scalar-minimizer input needs no uniform version. This gives the same positive upper/lower bounds on limiting slopes for the profile contradictions.

Child data, covariance margins (choose with room at the limit), and moving barriers in PE, Section~6 then use PE~(4) with that common \(T_m\), and the fixed finite depth subsequence proof applies unchanged. Thus the parameter replacement invokes no theorem uniform up to the critical temperature or change of the fixed propagation slack.

\paragraph{Scalar coercivity}
We first establish a coercive comparison. If \(K'=h_*+b z\) is an admissible scalar path with \(0\le z\le C\), then
\[
 f_{\rm spin}(K')+\tfrac14\int b z^2
 \ \ge F_*+c\int_0^{Q_0}\left(\int_s^{Q_0}(\alpha_{K'}-\alpha_t)\right)^2 ds ,
                                                                     \tag{G2}\label{eq:G2}
\]
where \(Q_0\) is a fixed clock bound, after annealed extension of
the time distributions. Use the base optimal feedback martingale
\(u_s=V_z(s,Z_s)\) as control in both scalar problems, as verified
in PE, Section~5. This gives the lower first variation
\(\frac12\int\Gamma(s)(\alpha_{K'}-\alpha_t)(s)ds\), together
with the terminal \(\log\cosh\) convexity remainder.

Write \(P(s)=\int_s^{Q_0}(\alpha_{K'}-\alpha_t)\) and
\(A=\int_0^{Q_0}P(s)^2ds\). The terminal displacement is
\(Y=\int(\alpha_{K'}-\alpha_t)u=\int P(s)\chi_s\,dB_s\),
where \(\chi_s=V_{zz}(s,Z_s)\). On bounded clocks,
\(\mathbf E\chi_s^2\ge c\), while \(\chi_s\) is bounded, so
\(\mathbf EY^2\ge cA\) and \(\mathbf EY^4\le CA^2\).
Also \(Y\) is deterministically bounded from its original integral.
Choose a fixed compact set for the base terminal argument with
complement probability at most \(c^2/(4C)\). Cauchy--Schwarz leaves
at least \(cA/2\) of \(\mathbf EY^2\) on this set. The second
derivative of \(\log\cosh\) is bounded below along every displaced
segment there. Its convexity remainder is therefore at least \(c'A\),
as required in \eqref{eq:G2}.

The first variation plus penalty change cancels linearly, leaving
\(\tfrac12\int dv\, b(z-q)^2\int_0^1(1-y)[1-b\Gamma'(t+y b(z-q))]dy\ge0\) by \eqref{eq:P5}, without needing monotonicity of \(z\). The lower bound on \(\mathbf E\chi_s^2\) here follows also by comparison to the base \(T\) scalar path on a fixed compact including its annealed extension as in \eqref{eq:P2}.

\paragraph{A constrained interpolation}
Choose \(e_0=x^{d_g}\), where the small exponent \(d_g>0\) will
be fixed below, and set \(d_{\rm PE}=e_0^{1/4}\). Apply PE with
\(d=d_{\rm PE}\), strength \(e_0\), and error parameter
\(E=x^{4-d_g/16}d_{\rm PE}=x^{4+3d_g/16}\). Taking its fixed
power slack sufficiently small relative to \(d_g\) satisfies
PE~(3) and makes \(m^{20\delta_{\rm PE}}E\le x^{4+c}\).
The endpoint scale is \(a_{\rm PE}\asymp x^{1-d_g/16}\gtrsim x\).
Also \(e_0/d_{\rm PE}^3\to0\), so the taper is identically 1:
this comparison removes the constant \(e_0\) and inserts \(e_0r\).
Equations~\eqref{eq:P2}--\eqref{eq:P4} and \eqref{eq:P9} verify
the endpoint and admissibility hypotheses with the common base
\(T_m\to T\) just discussed. The optimized value, including
\(\int e_0r^2/4\), is consequently within \(x^{4+c}\) of the
bare pressure.

We can pass to onto inverse paths \(\alpha:[0,2]\to[0,1]\) with constant bounds
\(\kappa=x^{4+1/4}\le\alpha'\le M_c=x^{-10}\), at upper-cost loss \(O(\kappa+1/M_c)\). Indeed any radius law on \([0,1]\) can be approximated in first transport moment to that accuracy, by a slight bounded-density spreading within the enlarged interval and mixing a uniform floor. Cost follows by \eqref{eq:A1}. Keep these latter constraints fixed and minimize
\[
 \Phi(\theta)=\min_r\{f_m(b-e,h_*+e r)+\textstyle\int e r^2/4\},\qquad
 e=(1-\theta)e_0+\theta b,\quad 0\le\theta\le1 .
\]
The density box is compact at fixed size, so a minimizer exists.
Equation~\eqref{eq:A1} gives monotonicity and Lipschitz continuity
of \(\Phi\), as well as its envelope upper derivative at a
minimizer. All paths are admissible. At the two endpoints,
\(\Phi(0)\le f_m(b,h_*)+x^{4+c}\) by the PE comparison and
\(\Phi(1)\ge F_*\) by \eqref{eq:G2}. At every intermediate
\(\theta\), trying a constrained approximation to \(q\) and
replacing the remaining matrix by the field \((b-e)q\) gives
\(\Phi(\theta)\le F_*+O(\kappa+M_c^{-1})\).

Fix \(\theta\) and a minimizing \(r\). The objective is an averaged
pressure, so this choice, its inverse \(\alpha\), and all intervals
chosen from them are deterministic before taking the positive-law
expectation. Evaluate \(S,B,D\) at this changed path. A subscript in
radius \(s\) means the cut at mass \(\alpha(s)\). Put
\[
 \widetilde b=b-e,\qquad
 K^{\rm emp}=h_*+er+\widetilde bS.
\]
This empirical scalar clock replaces each remaining matrix overlap by
its mean \(S\). It is nonnegative and nondecreasing.

It remains to show that \(\Phi\) increases by at most
\(O(x^{4+c})\). Its envelope derivative is bounded by the integral
of \(D^2+(S-r)^2\). We first use the minimizing density constraints
to control \(S-r\), then compare \(K^{\rm emp}\) to \(t\).
Deterministic radius buffers will give high moments of \(Q-S\);
the pair inverse will improve these to the required integrated bound
on \(D^2\). The final envelope estimate combines these four steps.

First variation gives a potential \(H\) with
\(H'(s)=e(\alpha(s))(s-S_{\alpha(s)})\). Indeed the variation against
\(\delta\alpha\) is \(-\int H'\delta\alpha/2\), by changing variables
in \eqref{eq:A1}. The inverse is bi-Lipschitz at fixed size, and the
required path moments are continuous under these perturbations by PE,
Section~2. After integration by parts this is a linear functional of
the density variation. Exchanging density between two locations, while
keeping total mass fixed, gives a constant shift of \(H\) for which
\(\alpha'=\kappa\) on \(H>0\) and \(\alpha'=M_c\) on \(H<0\).

A component of \(H<0\) has radius width at most \(M_c^{-1}\).
Its endpoint signs and monotonicity put \(S\) between its two radius
endpoints; at exterior boundaries use \(0\le S\le1\). On \(H=0\),
\(S=s\) almost everywhere. We call the complement of the positive
components \emph{nonfloor}. Thus \(|S-s|\le M_c^{-1}\) there, while
the floor has total mass at most \(2\kappa\). Integrating \(H'\) on
a negative component also gives \(H\ge-C(e_0+\theta)M_c^{-2}\).

We will use the following local consequence. Within distance \(O(h)\)
of a nonfloor point \(s_0\), where \(M_c^{-1}\ll h\),
\(S=s_0+O(h)\), provided one further interval of length \(h\) fits
on each required side. For if \(S\) exceeded this range, monotonicity
would keep \(s-S\) of one sign and of size \(\gtrsim h\) on such
an interval. Since \(e\asymp e_0+\theta\), integrating \(H'\) to
or from \(s_0\), where \(H(s_0)\le0\), would force a negative
value of size \(\gtrsim(e_0+\theta)h^2\). This contradicts the
lower bound for \(H\).

\paragraph{A quantitative coarse clock comparison}
The empirical scalar clock is a full mass path. Its time distribution
will be denoted by \(\alpha_{K^{\rm emp}}\), with field time as its
argument. We first need the coarse estimate
\[
 \|\alpha_{K^{\rm emp}}-\alpha_t\|_1\le x^{c_1},                         \tag{G3}\label{eq:G3}
\]
where \(c_1>0\) is independent of sufficiently small \(d_g\). Here and below these are scalar time profiles, not optimizer inverse paths. There is a monotone \(0\le z\le1\) with
\[
 f_{\rm spin}(h_*+e r+\widetilde b z)+\int\widetilde b z^2/4
       \le f_m(\widetilde b,h_*+e r)+x^{c_0}
\]
for some absolute \(c_0>0\). If \(\sup\widetilde b\lesssim x^{1/10}\),
take \(z=0\) and use the pressure Lipschitz bound in \eqref{eq:A1}.
Otherwise \(\widetilde b=(1-\theta)(b-e_0)\) has
\(b_*:=\inf\widetilde b\gtrsim x^{1/10}\), since \(b-e_0\asymp1\).
We use the quantitative bounds in the proof of PE, Section~5,
Lemma~5.7, as follows.

In that scalar-trial interpolation, denote its interpolation time by
\(\tau\), its inverse floor and cap by \(\eta_m=m^{-1/8}\) and
\(Q_m=m^{1/8}\), and its optimizer by \(z_\tau\). In this paragraph
the observables refer to this inner optimizer, with clocks
\(((1-\tau)\widetilde b,h_*+er+\tau\widetilde b z_\tau)\).
The mismatch
satisfies \(\int(S-z_\tau)^2dv\le\eta_m+Q_m^{-2}\). For
\(\tau\ge\varepsilon\), after also discarding masses below
\(\varepsilon\), the covariance trace satisfies
\[
 \int_\varepsilon^1\mathbf E\operatorname{tr}C_v^2dv
       \le R:=\frac{Q_m}{m b_*\varepsilon^3},
\]
and the variance of an unnormalized \(L\)-interval integral is at
most \(C Q_m/(m b_*\varepsilon^2)\). Partition the mass interval
where \(S\) varies by at most \(\Delta\), and discard end strips
of width \(\rho\). There are at most \(C(1+\Delta^{-1})\) strips.
The omitted masses and times therefore cost
\(O(\varepsilon+\rho(1+\Delta^{-1}))\). On every retained mass,
the earlier and later \(L\)-averages give
\[
 \bigl\||X_v|^2-S_v\bigr\|_2
 \le C\Delta+\frac C\rho
       \left(\frac{Q_m}{m b_*\varepsilon^2}\right)^{1/2}=:A.
\]
Finally use
\(D_v^2=\operatorname{Var}(|X_v|^2)+\mathbf E\operatorname{tr}C_v^2
+2\mathbf E X_v^{\mathsf T}C_vX_v\), with the last term bounded
in integrated mass by \(2R^{1/2}\). The integrated envelope derivative
is consequently at most
\[
 C\{\eta_m+Q_m^{-2}+\varepsilon+\rho(1+\Delta^{-1})
                   +A^2+R+R^{1/2}\}.
\]
With \(\varepsilon=\Delta=x^{1/100}\) and \(\rho=x^{2/100}\),
each term saves a fixed power of \(x\): for example,
\(R\lesssim x^{128/25}\) and
\(A\lesssim x^{1/100}+x^{509/200}\).
These estimates use only bounded monotone actual clocks and the lower
bound on \(b_*\), so they are uniform over our minimizing \(r\).
The scalar endpoint of this interpolation supplies the asserted trial.
We now return to the observables of the outer minimizing \(r\).

Put \(K'=h_*+er+\widetilde b z\) and
\(\bar z=(er+\widetilde b z)/b\). The exact penalty decomposition is
\[
 er^2+\widetilde b z^2
       =b\bar z^2+\frac{e\widetilde b}{b}(r-z)^2.
\]
Apply \eqref{eq:G2} to \(K'=h_*+b\bar z\). The trial inequality and
\(\Phi(\theta)\le F_*+O(\kappa+M_c^{-1})\) then bound both
the squared primitive in \eqref{eq:G2} and
\(\int e\widetilde b(r-z)^2dv\) by \(O(x^{c_0})\), after
decreasing \(c_0\).

To pass from the primitive to the profiles, write
\(g=\alpha_{K'}-\alpha_t\) and \(P(s)=\int_s^{Q_0}g\).
The total variation of \(g\) is bounded. Averaging over intervals
of length \(a\) gives
\(\|g\|_1\le C(a+a^{-1}\|P\|_1)\); boundary intervals have
the same \(O(a)\) cost. Since \(\|P\|_1\lesssim x^{c_0/2}\),
choosing \(a=x^{c_0/4}\) gives a power-small profile error.
Moreover the cap/floor mismatch and Cauchy--Schwarz give
\[
 \int\widetilde b|S-z|dv
 \lesssim\kappa+M_c^{-1}+e_0^{-1/2}x^{c_0/2}.
\]
For monotone clocks, the \(L^1\) distance between their quantiles
equals that between their time distributions. Thus this is also the
cost of replacing \(K'\) by \(K^{\rm emp}\). Choose \(d_g\)
sufficiently small relative to \(c_0\); a fixed smaller exponent
\(c_1>0\) now proves \eqref{eq:G3} independently of \(d_g\).
Decrease \(c_1\) once more, if needed, to cover the bounded-clock
scalar derivative comparison used at the end of the proof; all
remaining exponents are chosen after this choice.

\paragraph{Integrated pinned moments and rough concentration}
Choose \(\delta=x^{a_g}\) with \(0<a_g\ll c_1\), put
\(h=\delta^2\), and take \(d_g\) sufficiently small relative to
\(a_g\). Near a nonfloor point \(s_0\ge\delta\), the local
potential bound gives \(S=s_0+O(h)\), and hence
\(K^{\rm emp}\ge c\delta\). Such a clock cannot occur at mass
\(\ll\delta\): monotonicity and the regular initial slope of \(t\)
would give an area of order \(\delta^2\) between the two time
distributions, contradicting \eqref{eq:G3}, since \(c_1>2a_g\).
Thus these neighborhoods and all later cuts have mass
\(\gtrsim\delta\). They have room for the required buffers, since
\(s_0\le1+M_c^{-1}\), while the radius interval extends to 2.

For fixed \(p\ge2\), on a radius window \(I\) there (within \([0,5/4]\)) we have
\[
 \left\|\int_I W\,ds\right\|_{p/2}\lesssim_{p,\circ} x^6/(e_0\delta^2).   \tag{G4}\label{eq:G4}
\]
\paragraph{A reusable residual estimate.}
\label{par:global-block-estimates}
We isolate the deterministic blocking calculation also used in
Sections~\ref{sec:budgets} and~\ref{sec:statistics}. It applies to the full single-path
positive law and does not require an optimizing radius.

\begin{lemma}[Residual blocks and integrated increments]
\label{global:residual-blocks}
Fix a positive path law with terminal mass \(a\in[1,4]\), and
parametrize a deterministic increasing family of revealed cuts by \(s\). At a cut \(s\), let branches \(\mathsf a,\mathsf b\)
be conditionally independent continuations with their normalized positive
kernels, given the common prefix. Partition their futures into the same
\(k\) deterministic blocks, the last ending at the terminal spin.
Write \(\Delta_j\) for the martingale increments, so that
\(y-X_s=\sum_{j=1}^k\Delta_j\), and let \(\mu'_j\le a\)
be the largest mass in block \(j\). Evaluating \(W_s\) on branch
\(\mathsf a\), we have
\begin{equation}
 W_s\le k\sum_{i,j}\mu'_j\,
    \mathbf E_{\mathsf b}[(\Delta_i^{\mathsf a}\cdot\Delta_j^{\mathsf b})^2\mid\mathsf a].
 \label{eq:global-block-decomposition}
\end{equation}
Conditioning on \(\mathsf a\) retains its whole branch and the common prefix;
only continuation \(\mathsf b\) is averaged. For any deterministic interval
\(J\), any increment length \(l>0\), and \(q\ge1\), provided all
indicated cuts exist, put \(s_-=\inf J\) and \(s_+=\sup J+2l\). Then
\begin{equation}
 \left\|\int_J|X_{s+2l}-X_{s+l}|^2ds\right\|_q
       \le C_q l\,\bigl\||X_{s_+}-X_{s_-}|^2\bigr\|_q .
 \label{eq:global-grid}
\end{equation}
The same estimate holds with increment \(X_{s+l}-X_s\).
For any nonnegative process \(V_s\), not necessarily adapted,
\begin{equation}
 \left\|\int\mathbf E_sV_s\,ds\right\|_q
       \le C_q\left\|\int V_s\,ds\right\|_q .
 \label{eq:global-dual-doob}
\end{equation}
\end{lemma}

\begin{proof}[Proof of the lemma]
The step Hessian rule writes \(H_s\) as the conditional sum of future
increment covariances weighted by their masses; the last term is the
terminal Gibbs covariance multiplied by \(a\). Conditional orthogonality
gives
\[
 H_s\preceq\sum_j\mu'_j\mathbf E_s
                 \Delta_j\Delta_j^{\mathsf T}.
\]
Expanding only branch \(\mathsf a\)'s residual and applying Cauchy--Schwarz
gives \eqref{eq:global-block-decomposition}. In particular
\(H_s\preceq aC_s\). The fixed-size approximation in
\eqref{eq:A1} passes this statement to filled clocks.

For the integrated estimate, the shifted-grid identity is
\[
 \int_J|X_{s+2l}-X_{s+l}|^2ds
 =\int_0^l\sum_{j:\,u+jl\in J}
       |X_{u+(j+2)l}-X_{u+(j+1)l}|^2du.
\]
At each fixed \(u\) the increments are disjoint. Apply the martingale
square-function and Doob inequalities to \(X-X_{s_-}\), and integrate
in \(u\), to obtain \eqref{eq:global-grid}; for \(q=1\) martingale
orthogonality suffices. The earlier increment has the same proof.

For \eqref{eq:global-dual-doob} with \(q>1\), pair with a nonnegative
\(Z\) of unit \(L^{q'}\) norm. Move the conditional expectation onto
\(Z\), bound \(\mathbf E_sZ\) by its supremum, and apply Doob and
H\"older. For \(q=1\), the two expectations are equal.
\end{proof}

\paragraph{Application to the integrated pinned moment.}
We return to the normalized terminal mass \(a=1\) and prove
\eqref{eq:G4}. Here \(H_s\), the source Hessian, is distinct from
the optimizer potential \(H(s)\). The lemma also applies after an
increment on branch \(\mathsf b\) is averaged back to its fork, by
\eqref{eq:global-dual-doob}. We use \(\mu'_j\le1\), or directly
\(H_s\preceq C_s\), giving
\[
 W_s\le\mathbf E_{\mathsf b}[((y^{\mathsf a}-X_s)\cdot(y^{\mathsf b}-X_s))^2\mid\mathsf a].
\]

For the residual decomposition in \eqref{eq:G4}, fix a small constant
\(L_*>0\) with \(s+L_*<2\) for \(s\in I\). Choose
\(l_0=2^{-J_*}L_*\) of arbitrarily small inverse-polynomial size.
Use the initial block \(X_{s+l_0}-X_s\), the blocks
\(X_{s+2l}-X_{s+l}\) for \(l=l_0,2l_0,\ldots,L_*/2\), and
the terminal residual \(y-X_{s+L_*}\). There are
\(k=O(\log(1/x))\) blocks and their sum is exactly \(y-X_s\).

For any pair other than two initial blocks, project the block with
the later starting distance, denoted \(l_{\max}\). The anchor
\(A\) is the other branch's increment. Write
\(A=|A|\widehat A\), with \(\widehat A=0\) when \(A=0\),
and apply \eqref{eq:A3} to \(\widehat A\). Conditional on the
common prefix, the whole other branch is exterior data, so fixing
this anchor leaves the projected branch's normalized future kernel
unchanged. Equal starting distances may be resolved in either order.
The gap after the fork has field length at least \(e_0l_{\max}\)
and mass at least \(c\delta\). For a finite block, conditional
Jensen converts the residual estimate at its start into the same
estimate for its increment.

Fix a small \(\epsilon>0\). At a sufficiently high fixed even
moment, Markov's inequality gives
\[
 \mathbf P\left\{
 |\widehat A\cdot\Delta_{\rm later}|^2
       >\frac{x^{6-\epsilon}}{e_0\delta^2l_{\max}}\right\}
       =O(x^K)
\]
for any prescribed fixed \(K\). Indeed \eqref{eq:A3} gives the
same square scale with \(x^6\log^4N\); taking the moment order
large relative to \(K/\epsilon\) supplies this tail bound. The
factor \(|A|^2\) remains in the subsequent integral. On the
exceptional event the unnormalized squared product is bounded,
since \(|y|,|X|\le1\). Jensen before and after conditional averaging
shows that its integral has \(L^{p/2}\) norm \(O(x^{2K/p})\).
Choose \(K\) large enough to make this negligible.

By \eqref{eq:global-grid}--\eqref{eq:global-dual-doob}, the earlier
squared anchor integrated over \(I\) has \(L^{p/2}\) norm at most
\(C_p l_{\max}\). If both blocks are terminal, boundedness gives
the same assertion since \(l_{\max}=L_*\) is fixed. Thus each
projected pair contributes at most
\(C_p x^{6-\epsilon}/(e_0\delta^2)\). For two initial blocks,
bound one by a constant and integrate the other's square; the cost
is \(O_p(l_0)\), negligible by the choice of \(l_0\).
Summing the logarithmically many pairs in
\eqref{eq:global-block-decomposition} proves \eqref{eq:G4}.

\paragraph{Earlier and later buffers.}\label{par:global-brackets}
We next turn the integrated pinned estimate into control at individual
cuts. The following bracket argument also records what is needed when
the later applications use different radius windows.

\begin{lemma}[Buffer brackets and a common past]
\label{global:buffer-brackets}
Fix a full single-path positive law with terminal mass \(a\in[1,4]\)
and deterministic revealed cuts. Let \(w\) be one such cut, and let
\(\rho_-\), \(\rho_+\) be deterministic probability measures supported
on cuts before and after \(w\), respectively. Write
\[
 A_-=\int L_v\,\rho_-(dv),\qquad
 A_+=\int L_v\,\rho_+(dv).
\]
Suppose that, for a deterministic number \(\bar S\),
\[
 |S_v-\bar S|\le\eta\quad
  \text{on the two buffers},\qquad
 \|A_\pm-\mathbf EA_\pm\|_p\le\varepsilon_p,
 \qquad 2\le p<\infty.
\]
The bound on \(S_v\) is required only almost everywhere for the two
buffer measures. Then
\[
 \bigl\||X_w|^2-\bar S\bigr\|_p+
 \left\|\int|X_w-X_v|^2\rho_-(dv)\right\|_p
       \le C(\eta+\varepsilon_p).
\]
If two cuts \(w,w'\) use the same earlier measure \(\rho_-\), supported
before both, and each has a later buffer satisfying the same hypotheses
with the same \(\bar S,\eta,\varepsilon_p\), then
\[
 \bigl\||X_w-X_{w'}|^2\bigr\|_p\le C(\eta+\varepsilon_p).
\]
\end{lemma}

\begin{proof}[Proof of the lemma]
Since \(2\mathbf EL_v=S_v\), the buffer assumptions give
\(|2\mathbf EA_\pm-\bar S|\le\eta\). The martingale identities are
\[
 \begin{split}
 2\mathbf E_w A_-&=|X_w|^2-
           \int|X_w-X_v|^2\rho_-(dv),\\
 2\mathbf E_w A_+&=\int\mathbf E_w|X_v|^2\rho_+(dv)
                      \ge |X_w|^2.
 \end{split}
\]
Thus \(|X_w|^2\) lies between \(2\mathbf E_wA_-\) and
\(2\mathbf E_wA_+\), each of which is within
\(\eta+2\varepsilon_p\) of \(\bar S\) in \(L^p\), by conditional
Jensen. The first identity then also bounds its nonnegative past
squared-increment integral. For two cuts, integrate
\[
 |X_w-X_{w'}|^2\le
       2|X_w-X_v|^2+2|X_{w'}-X_v|^2
\]
against their common earlier measure and use the two preceding bounds.
\end{proof}

Apply the lemma with normalized radius averages over buffers of length
comparable to \(h\), before and after cuts near \(s_0\). Their measures
are deterministic because the optimizer was fixed before sampling.
For an interval of radius length \(h\), the mass weight in
\eqref{eq:A1} is \(k(v)=1/[h\alpha'(r(v))]\). Consequently
\eqref{eq:G4} gives
\[
 \|A_\pm-\mathbf EA_\pm\|_p
 \lesssim_{p,\circ}\frac{x^3}{\sqrt{e_0}\delta h\sqrt\kappa}
 =x^{7/8-d_g/2-3a_g}\,O_{p,\circ}(1)=O_p(h).
\]
The last estimate holds by taking the fixed exponents small before
choosing the subpower loss. The local potential bound gives
\(S_v=s_0+O(h)\) throughout these buffers. Lemma~\ref{global:buffer-brackets}
therefore applies with \(\bar S=s_0\), \(\eta=O(h)\), and
\(\varepsilon_p=O_p(h)\). It bounds both
\(\||X_w|^2-s_0\|_p\) and the averaged past squared increments by
\(O_p(h)\). Using the same earlier buffer for two nearby cuts gives
\(\||X_w-X_{w'}|^2\|_p=O_p(h)\).

Now take two branches splitting at \(s_0\), and project their
terminal spins to cuts a distance comparable to \(h\) later.
Their squared norms are \(s_0+O_{L^p}(h)\), and their increments
from the common ancestor have squared norms \(O_{L^p}(h)\).
Polarization therefore puts their projected inner product within
\(O_{L^p}(h)\) of \(s_0\). The two projection errors in
\eqref{eq:A3} are at most
\(O_{p,\circ}(x^3/(\delta\sqrt{e_0h}))\), so in particular
\(\|Q-S\|_p=O_p(\delta)\) at this nonfloor split.

For splits below radius \(\delta\), choose instead a nonfloor cut
at radius comparable to a sufficiently large fixed multiple of
\(\delta\). To find it, take a mass band of that scale.
Equation~\eqref{eq:G3} and monotonicity put \(K^{\rm emp}\)
between comparable multiples of \(\delta\) on an interior band;
the discarded floor has mass at most \(2\kappa\ll\delta\).
At a nonfloor point of this band, \(h_*=o(\delta)\) and
\(r=S+O(M_c^{-1})\), so
\(K^{\rm emp}=br+o(\delta)\) and \(r\asymp\delta\), as needed.
The same buffers give \(\||X|^2\|_p=O_p(\delta)\) there.
Project both branches to this cut, using its last comparable
subinterval for \eqref{eq:A3}; Cauchy--Schwarz bounds their projected
product by \(O_p(\delta)\). Since \(S\lesssim\delta\) at the
original cut, we obtain, after weakening the exponent,
\(\|Q-S\|_p\lesssim_p x^{a_g/2}\) at almost every nonfloor split,
for every fixed moment order.

\paragraph{Centered inversion and absorption}
At the sampled clock \(K^{\rm emp}\), the matrix multiplier
\(\widetilde b\) has the gap required in \eqref{eq:H1}, of size
at least \(ce_0\). To check this, \eqref{eq:G3} and the bounded-clock
PDE estimates compare scalar derivatives at a common time with error
\(O(x^{c_1})\).
For \(v\ge\delta\), monotonicity and the same area argument give
\(K^{\rm emp}(v)\gtrsim v\). The positive part of
\(b(v)-b_\circ\) is at most \(Cc_b\lambda v^2\), which is
absorbed by \eqref{eq:P2} when \(c_b\) is small. Below \(\delta\)
this positive part is \(o(e_0)\). Finally, replacing \(b\) by
\(b-e\) gains \(e\Gamma'\): wherever this is not comparable to
\(e_0\), \(\Gamma'\) is small and the original gap is bounded
below by a constant. The common-time error is \(o(e_0)\), so the
gap persists in both cases.

Fix the old pair \(ij\) and use \eqref{eq:A7} to first order about
\(K^{\rm emp}\), with test \(T_{ij}=Q_{ij}-S_{ij}\).
The row mismatch is exactly
\[
 J=(h_*+er)+\widetilde b Q-K^{\rm emp}
   =\widetilde b(Q-S).
\]
The test and the deterministic centers \(S\) remain fixed throughout
this cavity interpolation; the actual row and the independent scalar
row at \(K^{\rm emp}\) are its positive-semidefinite endpoints.
The unknown pair function is the deterministic, topology-indexed
expectation
\(u_{ab}=\nu[T_{ij}(Q_{ab}-S_{ab})]\).
It is unchanged on adding unused labels, by positive-law
marginalization. The scalar degree-zero term, including any
deterministic offset on a given input topology, vanishes because
\(\nu[T_{ij}]=0\). The linear terms give \eqref{eq:H1} with
multiplier \(\widetilde b\). Its errors are \(O(x^6)\) and
products of \(T_{ij}\) with at least two \(Q-S\) insertions.

Apply the inverse, evaluate at \(ab=ij\), and then integrate the
old fork mass. Absolute attachment domination bounds new fork
densities by a polynomial in \(e_0^{-1}\); a fork reused at \(ij\)
is now integrated as well. Truncate one error insertion at
\(x^{a_g/4}\). Its higher fixed moments, just proved off the floor,
make the exceptional probability \(O(x^6)\) by Markov's inequality,
and all factors are bounded there. To charge a floor placement,
follow its fork to the first time that vertex was introduced. It is
either the original \(ij\) fork, whose mass is now integrated, or a
fresh fork with bounded density. Later reuses retain the ordinary
bounded atom weights and introduce no new fixed mass. The floor mass
is \(O(\kappa)\), so the finite union of these placements also costs
\(O(\kappa)\) under the dominated allocation measure.
On the retained event this leaves \(x^{a_g/4}|T_{ij}(Q_{ab}-S_{ab})|\).
Cauchy--Schwarz bounds its expectation by
\(x^{a_g/4}D_{v_{ij}}D_{v_{ab}}\); the bounded fork measures and
\(2D_vD_w\le D_v^2+D_w^2\) give the integrated estimate
\[
 \int D_v^2dv\le O(e_0^{-C})\big(\kappa+x^6+x^{a_g/4}\int D_v^2dv\big).
\]

For nonstep paths, perform this inversion first at fixed \(m\) on
step approximations, using each approximation's own \(S\) and
\(K^{\rm emp}\) in the cavity formula. Their moments converge at
almost every placement and in mass \(L^1\). The common-time scalar
comparison and monotone lower clock bound persist, so the sheet gap
remains positive when \(b,e\) are sampled from their paths.
Pass to the limit inside the dominated integrated error bounds, as
in the \hyperref[par:inverse-limits]{fixed-resolution limit argument}
of Proposition~\ref{inv:pair-inverse}.

Now choose \(d_g\) so small that \(Cd_g<a_g/4\) and
\(Cd_g<1/4\). The first inequality absorbs the variance term on
the right; the second makes \(e_0^{-C}\kappa\) smaller than
\(x^4\) by a fixed power. At a minimizing \(r\), \eqref{eq:A1}
gives the envelope upper derivative
\[
 \Phi'(\theta)\le\frac14\int(b-e_0)
                  \{D_v^2+(S_v-r(v))^2\}\,dv
                  =O(x^{4+c}).
\]
Here the mismatch integral is \(O(\kappa+M_c^{-2})\), by the
floor mass and the nonfloor cap bound. Integrate over
\(0\le\theta\le1\), and use
\(F_*\le\Phi(1)\) and \(\Phi(0)\le f_m^0+x^{4+c}\).
Decreasing the fixed saving exponent if necessary proves \eqref{eq:G1}.

\end{proof}
